%Data institusi

%\input{Dozentdaten}

%\renewcommand{\fachbereich}{Bidang studi}

%\renewcommand{\dozent}{Prof. Dr. . }

%Data ujian

\renewcommand{\klausurgebiet}{ }

\renewcommand{\klausurtyp}{ }

\renewcommand{\klausurdatum}{ . 20}

\klausurvorspann {\fachbereich} {\klausurdatum} {\dozent} {\klausurgebiet} {\klausurtyp}

%Data untuk tabel poin berikut

\renewcommand{\aeins}{  3  }

\renewcommand{\azwei}{  3   }

\renewcommand{\adrei}{  6  }

\renewcommand{\avier}{  3  }

\renewcommand{\afuenf}{  3  }

\renewcommand{\asechs}{  5  }

\renewcommand{\asieben}{  13  }

\renewcommand{\aacht}{  4  }

\renewcommand{\aneun}{  3  }

\renewcommand{\azehn}{  7  }

\renewcommand{\aelf}{  6  }

\renewcommand{\azwoelf}{  3  }

\renewcommand{\adreizehn}{  6   }

\renewcommand{\avierzehn}{  65  }

\renewcommand{\afuenfzehn}{    }

\renewcommand{\asechzehn}{    }

\renewcommand{\asiebzehn}{    }

\renewcommand{\aachtzehn}{    }

\renewcommand{\aneunzehn}{    }

\renewcommand{\azwanzig}{    }

\renewcommand{\aeinundzwanzig}{    }

\renewcommand{\azweiundzwanzig}{    }

\renewcommand{\adreiundzwanzig}{    }

\renewcommand{\avierundzwanzig}{    }

\renewcommand{\afuenfundzwanzig}{    }

\renewcommand{\asechsundzwanzig}{    }

\punktetabelledreizehn

\klausurnote

\newpage

\setcounter{section}{0}



\inputaufgabeklausurloesung
{3}

{

Definisikan istilah-istilah berikut
\zusatzklammer {yang dicetak miring} {} {.}
\aufzaehlungsechs{Suatu
\stichwort {ruang Hausdorff} {}
$X$.

}{Suatu
\stichwort {pemetaan diferensiabel} {}
\maabbdisp {\varphi} { L } { M
} {}
antara dua manifold diferensiabel
\mathkor {} {L} {dan} {M} {.}

}{
\stichwort {Bundel tangen} {}
dari suatu manifold diferensiabel $M$.

}{
\stichwort {Matriks fundamental kedua} {}
dari parametrisasi lokal
\maabbdisp {} {V} {U \subseteq Y
} {}
\zusatzklammer {dengan

\mavergleichskette
{\vergleichskette
{ V
}
{ \subseteq }{ \R^{n-1}
}
{  }{
}
{    }{
}
{   }{
}
}
{}{}{}
terbuka} {} {}
suatu hiperpermukaan diferensiabel

\mavergleichskette
{\vergleichskette
{ Y
}
{ \subseteq }{ \R^n
}
{  }{
}
{    }{
}
{   }{
}
}
{}{}{.}

}{
\stichwort {Turunan eksterior} {} dari suatu bentuk diferensial terdiferensialkan secara kontinu
\mathl{\omega \in
{ \mathcal E }^{ k } (  M   )}{} pada manifold diferensiabel $M$.

}{Suatu \stichwort {partisi kesatuan yang tersubordinasi pada penutup} {,} dengan

\mavergleichskette
{\vergleichskette
{  X
}
{ = }{ \bigcup_{i \in I} W_i
}
{  }{
}
{    }{
}
{   }{
}
}
{}{}{}
suatu
\definitionsverweis {penutup terbuka}{}{}
dari
\definitionsverweis {ruang topologis}{}{}
$X$.
}

}
{

\aufzaehlungsechs{Suatu
\definitionsverweis {ruang topologis}{}{}
$X$ disebut Hausdorff jika untuk setiap dua titik berbeda
\mathl{x, y \in X}{} terdapat dua
\definitionsverweis {himpunan terbuka}{}{}
\mathkor {} {U} {dan} {V} {}
yang memenuhi
\mathl{x \in U, \, y \in V}{} dan
\mathl{U \cap V = \emptyset}{.}
}{Misalkan
\mathkor {} {(U_i,U_i', \alpha_i, i \in I)} {dan} {(V_j,V_j', \beta_j, j \in J)} {}
adalah
\definitionsverweis {atlas}{}{}
dari
\mathkor {} {M} {dan} {N} {.}
Pemetaan
\maabbdisp {\varphi} { L } { M
} {}
disebut diferensiabel jika kontinu dan jika untuk setiap
\mathkor {} {i \in I} {dan setiap} {j \in J} {}
pemetaan-pemetaan
\maabbdisp {\beta_j \circ  \varphi \circ (\alpha_i)^{-1}} {\alpha_i( \varphi^{-1}(V_j) \cap U_i) } { V_j'
} {}
\definitionsverweis {terdiferensialkan secara kontinu}{}{.}
}{Himpunan

\mavergleichskettedisp
{\vergleichskette
{TM
}
{ =} {\biguplus_{P \in M} T_PM
}
{ } {
}
{ } {
}
{ } {
}
}
{}{}{,}
yang dilengkapi dengan pemetaan proyeksi
\maabbeledisp {\pi} {TM} {M
} {(P,v)} {P
} {,}
disebut bundel tangen dari $M$.
}{Misalkan $Y$
\definitionsverweis {diorientasikan}{}{}
oleh suatu
\definitionsverweis {medan normal satuan}{}{}
$N$, dengan $N$ dipandang sebagai medan pada $V$. Definisikan

\mavergleichskettedisp
{\vergleichskette
{ h_{i j}
}
{ =} {  \left\langle  \partial_i \partial_j \varphi , N  \right\rangle
}
{ } {
}
{ } {
}
{ } {
}
}
{}{}{.}
Matriks fundamental kedua
dari $\varphi$ adalah matriks
\zusatzklammer {yang bergantung pada

\mavergleichskettek
{\vergleichskettek
{ Q
}
{ \in }{ V
}
{  }{
}
{    }{
}
{   }{
}
}
{}{}{}} {} {}

\mavergleichskettedisp
{\vergleichskette
{ H
}
{ =} {  \left(  h_{ij}   \right)_{1 \leq i,j \leq n-1}
}
{ } {
}
{ } {
}
{ } {
}
}
{}{}{.}
}{Turunan eksterior dari $\omega$, secara lokal pada suatu peta tempat $\omega$ berbentuk

\mavergleichskettedisp
{\vergleichskette
{ \omega
}
{ =} {\sum_{ { \# \left(   I  \right) } = k , \, I \subseteq \{1 , \ldots ,  \operatorname{dim}\, M \}  }f_I dx_I
}
{ } {
}
{ } {
}
{ } {
}
}
{}{}{}
didefinisikan oleh

\mavergleichskettedisp
{\vergleichskette
{  d \omega
}
{ =} {\sum_{ { \# \left(   I  \right) } = k   } d(f_I ) \wedge dx_I
}
{ } {
}
{ } {
}
{ } {
}
}
{}{}{}
rumus di atas.
}{Suatu keluarga fungsi
\maabbdisp {h_j} {X} {\R
} {}
dengan

\mavergleichskette
{\vergleichskette
{ j
}
{ \in }{ J
}
{  }{
}
{    }{
}
{   }{
}
}
{}{}{}
disebut partisi kesatuan yang tersubordinasi pada penutup jika sifat-sifat berikut berlaku.
\aufzaehlungvier{Berlaku

\mavergleichskette
{\vergleichskette
{ h_j(X)
}
{ \subseteq }{ [0,1]
}
{  }{
}
{    }{
}
{   }{
}
}
{}{}{}
untuk setiap

\mavergleichskette
{\vergleichskette
{ j
}
{ \in }{ J
}
{  }{
}
{    }{
}
{   }{
}
}
{}{}{.}
}{Setiap titik

\mavergleichskette
{\vergleichskette
{ P
}
{ \in }{ X
}
{  }{
}
{    }{
}
{   }{
}
}
{}{}{}
memiliki suatu
\definitionsverweis {lingkungan terbuka}{}{}

\mavergleichskette
{\vergleichskette
{ P
}
{ \in }{  U
}
{  }{
}
{    }{
}
{   }{
}
}
{}{}{}
sedemikian sehingga
\definitionsverweis {fungsi-fungsi yang dibatasi}{}{}

\mathl{h_j  \vert_{ U }}{} merupakan fungsi nol, kecuali untuk sejumlah berhingga indeks.
}{Berlaku

\mavergleichskette
{\vergleichskette
{  \sum_{j \in J} h_j
}
{ = }{ 1
}
{  }{
}
{    }{
}
{   }{
}
}
{}{}{.}
}{Untuk setiap

\mavergleichskette
{\vergleichskette
{ j
}
{ \in }{ J
}
{  }{
}
{    }{
}
{   }{
}
}
{}{}{}
terdapat suatu himpunan terbuka
\mathl{W_{i(j)}}{} dari penutup sedemikian sehingga
\definitionsverweis {tumpuan}{}{}
dari $h_j$ termuat dalam
\mathl{W_{i(j)}}{.}
}
}

 }



\inputaufgabeklausurloesung
{3}

{

Nyatakan teorema-teorema berikut.
\aufzaehlungdrei{
\stichwort {Teorema tentang sifat aljabar pemetaan Weingarten} {.}}{Teorema tentang keterorientasian dan bentuk volume positif.}{
\stichwort {Teorema titik tetap Brouwer} {.}}

}
{

\aufzaehlungdrei{\faktsituation {Misalkan

\mavergleichskette
{\vergleichskette
{ W
}
{ \subseteq }{ \R^n
}
{  }{
}
{    }{
}
{   }{
}
}
{}{}{}
\definitionsverweis {terbuka}{}{,}
\maabb {h} { W } { \R
} {}
adalah suatu
\definitionsverweis {fungsi yang dua kali terdiferensialkan secara kontinu}{}{}
dan

\mavergleichskette
{\vergleichskette
{ Y
}
{ = }{ h^{-1}(c)
}
{  }{
}
{    }{
}
{   }{
}
}
{}{}{}
adalah serat di atas

\mavergleichskette
{\vergleichskette
{ c
}
{ \in }{ \R
}
{  }{
}
{    }{
}
{   }{
}
}
{}{}{,}
dengan $h$
\definitionsverweis {reguler}{}{}
di setiap titik $Y$.}
\faktfolgerung {Maka
\definitionsverweis {pemetaan Weingarten}{}{}
$L_P$ di setiap titik

\mavergleichskette
{\vergleichskette
{ P
}
{ \in }{ Y
}
{  }{
}
{    }{
}
{   }{
}
}
{}{}{}
bersifat
\definitionsverweis {adjoin-diri}{}{.}}
\faktzusatz {}
\faktzusatz {}}{\faktsituation {Misalkan $M$ adalah suatu
\definitionsverweis {manifold diferensiabel
}{}{}
dengan suatu
\definitionsverweis {basis topologi terhitung}{}{.}}
\faktfolgerung {Terdapat suatu
\definitionsverweis {bentuk volume}{}{}
\definitionsverweis {kontinu}{}{}
tanpa titik nol pada $M$ jika dan hanya jika $M$
\definitionsverweis {dapat diorientasikan}{}{.}}
\faktzusatz {Bentuk volume ini juga dapat dipilih terdiferensialkan secara kontinu dan positif.}
\faktzusatz {}}{Misalkan
\maabbdisp {\psi} { B  \left( 0,r \right) } { B  \left( 0,r \right)
} {}
adalah suatu pemetaan yang terdiferensialkan secara kontinu dari bola tertutup di $\R^n$ ke dirinya sendiri. Maka $\psi$ memiliki suatu titik tetap.}

 }



\inputaufgabeklausurloesung
{6}

{

Tentukan
\definitionsverweis {kelengkungan}{}{}
di setiap titik lingkaran satuan yang diberikan secara implisit oleh

\mavergleichskettedisp
{\vergleichskette
{ h(x,y)
}
{ =} { x^2+y^2
}
{ =} { 1
}
{ } {
}
{ } {
}
}
{}{}{}
dengan
Lemma 3.11 (Geometri Diferensial (Osnabrück 2023))
dan medan gradien $h$.

}
{

\definitionsverweis {Medan gradien}{}{}
dari $h$ adalah $\begin{pmatrix}  2x \\ 2y   \end{pmatrix}$, sehingga suatu medan normal satuan adalah

\mavergleichskettedisp
{\vergleichskette
{ N \begin{pmatrix}  x  \\ y   \end{pmatrix}
}
{ =} { { \frac{   1 }{  \sqrt{x^2+y^2}  } } \begin{pmatrix}  x  \\ y   \end{pmatrix}
}
{ } {
}
{ } {
}
{ } {
}
}
{}{}{.}
\definitionsverweis {Diferensial total}{}{}
dari $N$ di suatu titik

\mavergleichskettedisp
{\vergleichskette
{ P
}
{ =} { \begin{pmatrix}  x  \\ y   \end{pmatrix}
}
{ } {
}
{ } {
}
{ } {
}
}
{}{}{}
adalah

\mavergleichskettedisp
{\vergleichskette
{  \left(D N \right)_{P}
}
{ =} { \begin{pmatrix}  { \frac{   y^2 }{   \left(  x^2+y^2  \right)^{3/2}  } }   &   - { \frac{   xy }{   \left(  x^2+y^2  \right)^{3/2}  } }   \\    - { \frac{   xy }{   \left(  x^2+y^2  \right)^{3/2}  } }    &  { \frac{   x^2 }{   \left(  x^2+y^2  \right)^{3/2}  } }  \end{pmatrix}
}
{ } {
}
{ } {
}
{ } {
}
}
{}{}{.}
Jika diterapkan pada vektor singgung
\mathl{\begin{pmatrix}  y  \\ -x   \end{pmatrix}}{}, diperoleh

\mavergleichskettealign
{\vergleichskettealign
{ \left(D N \right)_{P} \begin{pmatrix}  y  \\ -x   \end{pmatrix}
}
{ =} { \begin{pmatrix}  y^2    &   - xy    \\  - xy   &  x^2  \end{pmatrix} \begin{pmatrix}  y  \\ -x   \end{pmatrix}
}
{ =} { \begin{pmatrix}  y \left(  y^2 +x^2  \right)  \\ -x \left(  y^2 +x^2  \right)    \end{pmatrix}
}
{ =} { \begin{pmatrix}  y   \\ -x    \end{pmatrix}
}
{ } {
}
}
{}
{}{.}
Oleh karena itu, menurut
Lemma 3.11 (Differentialgeometrie (Osnabrück 2023))
kelengkungan kurva di $P$ sama dengan $-1$.
\zusatzklammer {Hasil ini sesuai karena
\mathl{\begin{pmatrix}  y  \\ -x   \end{pmatrix}}{} dan medan gradien merepresentasikan orientasi standar, sedangkan
\mathl{\begin{pmatrix}  y  \\ -x   \end{pmatrix}}{} adalah arah singgung untuk gerak searah jarum jam.} {} {}

 }



\inputaufgabeklausurloesung
{3}

{

Misalkan $X$ suatu
\definitionsverweis {ruang Hausdorff}{}{.}
Tunjukkan bahwa
\definitionsverweis {diagonal}{}{}

\mavergleichskettedisp
{\vergleichskette
{ \triangle
}
{ =} { \left\{ (x,y) \in X \times X  \mid   x = y        \right\}
}
{ } {
}
{ } {
}
{ } {
}
}
{}{}{}
merupakan
\definitionsverweis {subhimpunan tertutup}{}{}
dalam
\definitionsverweis {ruang produk}{}{}

\mathl{X \times X}{}.

}
{

Kita harus menunjukkan bahwa komplemen

\mavergleichskettedisp
{\vergleichskette
{W
}
{ =} { X \times X \setminus \triangle
}
{ =} { \left\{  (x,y)   \mid   x,y \in X , \,  x \neq y       \right\}
}
{ } {
}
{ } {
}
}
{}{}{}
terbuka. Jadi, misalkan
\mathl{(x,y)}{} adalah suatu pasangan dengan

\mavergleichskette
{\vergleichskette
{ x
}
{ \neq }{ y
}
{  }{
}
{    }{
}
{   }{
}
}
{}{}{.}
Berdasarkan sifat Hausdorff yang diasumsikan, terdapat himpunan-himpunan terbuka saling lepas
\mathl{U,V \subseteq X}{} dengan
\mathl{x \in U}{} dan
\mathl{y \in V}{.} Berlaku
\mathl{(x,y) \in U \times V}{,} dan menurut definisi topologi produk,
\mathl{U \times V}{} adalah suatu himpunan terbuka dalam
\mathl{X \times X}{.} Karena keduanya saling lepas, dari
\mathl{(z,w) \in U\times V}{} langsung diperoleh

\mavergleichskette
{\vergleichskette
{z
}
{ \neq }{w
}
{  }{
}
{    }{
}
{   }{
}
}
{}{}{.}
Jadi,

\mavergleichskettedisp
{\vergleichskette
{(x,y)
}
{ \in} { U \times V
}
{ \subseteq} {  W
}
{ } {
}
{ } {
}
}
{}{}{}
dan $W$ merupakan gabungan himpunan-himpunan produk terbuka seperti itu, sehingga himpunan tersebut sendiri terbuka.

 }



\inputaufgabeklausurloesung
{3}

{

Tunjukkan bahwa ketiga manifold satu dimensi
\mathlistdisp {S^1 = \left\{ (x,y) \in \R^2  \mid   \betrag { (x,y) } = 1         \right\}} {} {\R} {dan} {\R \setminus \{0\}} {}
tidak ada dua di antaranya yang homeomorfik.

}
{

Sebagai himpunan bagian tertutup dan terbatas dari $\R^2$, lingkaran satuan bersifat kompak
\zusatzklammer {Satz von Heine-Borel} {} {.} Sebaliknya, garis real $\R$ dan $\R \setminus \{0\}$ tidak kompak karena keduanya merupakan himpunan bagian tak terbatas dari $\R$. Jadi,
\mathl{S^1 \not \cong \R,\, \R \setminus \{0\}}{.}

Garis real terhubung, sebagaimana mengikuti dari
Zwischenwertsatz
. Sebaliknya, $\R \setminus \{0\}$ tidak terhubung karena dapat ditulis

\mavergleichskettedisp
{\vergleichskette
{\R_+
}
{ =} {(\R \setminus \{0\}) \cap [0, \infty]
}
{ =} {(\R \setminus \{0\}) \cap \, ]0, \infty]
}
{ } {
}
{ } {
}
}
{}{}{}
sehingga diperoleh suatu himpunan bagian tak kosong yang sekaligus terbuka dan tertutup. Jadi,
\mathl{\R \,\not \cong \R \setminus \{0\}}{.}

 }



\inputaufgabeklausurloesung
{5}

{

Misalkan

\mavergleichskette
{\vergleichskette
{ G
}
{ \subseteq }{ \R^m
}
{  }{
}
{    }{
}
{   }{
}
}
{}{}{}
\definitionsverweis {terbuka}{}{}
dan misalkan

\mavergleichskette
{\vergleichskette
{ M
}
{ \subseteq }{ G
}
{  }{
}
{    }{
}
{   }{
}
}
{}{}{}
suatu
$n$-\definitionsverweis {dimensional}{}{}
\definitionsverweis {submanifold tertutup}{}{,}
yang
\definitionsverweis {berorientasi}{}{}
dan dilengkapi dengan struktur Riemann terinduksi serta
\definitionsverweis {bentuk volume kanonik}{}{}
$\omega$. Misalkan

\mavergleichskette
{\vergleichskette
{ W
}
{ \subseteq }{ \R^n
}
{  }{
}
{    }{
}
{   }{
}
}
{}{}{}
terbuka dan misalkan
\maabbdisp {\varphi} { W } { U
} {}
suatu
\definitionsverweis {difeomorfisme}{}{}
dengan himpunan terbuka

\mavergleichskette
{\vergleichskette
{ U
}
{ \subseteq }{ M
}
{  }{
}
{    }{
}
{   }{
}
}
{}{}.
Tunjukkan bahwa pada $W$ berlaku rumus

\mavergleichskettealignhandlinks
{\vergleichskettealignhandlinks
{\varphi^* \left(  \omega \vert_U  \right)
}
{ =} { \left( \det  \left( \left\langle  \begin{pmatrix}  { \frac{   \partial  \varphi_1   }{  \partial   x_i   } }  \\ \vdots \\  { \frac{   \partial  \varphi_m   }{  \partial   x_i   } }   \end{pmatrix}  ,  \begin{pmatrix}  { \frac{   \partial  \varphi_1   }{  \partial   x_j   } }  \\ \vdots \\  { \frac{   \partial  \varphi_m   }{  \partial   x_j   } }   \end{pmatrix}   \right\rangle \right)_{1 \leq i,j \leq n} \right)^{1/2}  dx_1 \wedge \ldots \wedge dx_n
}
{ =} { \left( \det  \left(  { \frac{   \partial  \varphi_1   }{  \partial   x_i   } } { \frac{   \partial  \varphi_1   }{  \partial   x_j   } }  + \cdots +  { \frac{   \partial  \varphi_m   }{  \partial   x_i   } } { \frac{   \partial  \varphi_m   }{  \partial   x_j   } }  \right)_{1 \leq i,j \leq n} \right)^{1/2}  dx_1 \wedge \ldots \wedge dx_n
}
{ } {
}
{ } {
}
}
{}
{}{}{.}

}
{

Persamaan kedua langsung diperoleh dengan mengevaluasi hasil kali dalam standar pada $\R^m$. Menurut definisi
\definitionsverweis {matriks fundamental metrik}{}{,}
untuk

\mavergleichskette
{\vergleichskette
{ Q
}
{ \in }{ W
}
{  }{
}
{    }{
}
{   }{
}
}
{}{}{}

\mavergleichskettealign
{\vergleichskettealign
{ g_{ij} (Q)
}
{ =} { \left\langle  T_Q(\varphi)(e_i)  ,  T_Q(\varphi)(e_j)   \right\rangle_{\varphi(Q)}
}
{ =} { \left\langle  T_Q(\varphi)(e_i)  ,  T_Q(\varphi)(e_j)   \right\rangle
}
{ =} { \left\langle  \left(  D \varphi  \right)_{ Q } \left(   e_i   \right)  ,  \left(  D\varphi  \right)_{ Q } \left(   e_j   \right)   \right\rangle
}
{ =} { \left\langle  \begin{pmatrix}  { \frac{   \partial  \varphi_1   }{  \partial   x_i   } } ( Q)  \\ \vdots \\  { \frac{   \partial  \varphi_m   }{  \partial   x_i   } } ( Q)   \end{pmatrix}  ,  \begin{pmatrix}  { \frac{   \partial   \varphi_1   }{  \partial   x_j   } } ( Q )  \\ \vdots \\  { \frac{   \partial  \varphi_m   }{  \partial   x_j   } } ( Q)   \end{pmatrix}   \right\rangle
}
}
{}
{}{,}
karena ruang singgung
\mathl{T_{\varphi(Q)}M}{} dilengkapi dengan hasil kali dalam yang diinduksi dari $\R^m$, pemetaan singgung dalam kasus lokal adalah diferensial total, dan entri-entrinya dapat dinyatakan dengan turunan parsial. Oleh karena itu, pernyataan tersebut mengikuti dari
Lemma 16.7 (Differentialgeometrie (Osnabrück 2023)).

 }



\inputaufgabeklausurloesung
{weiter}

{

Kita tinjau

\mavergleichskette
{\vergleichskette
{ \R^6
}
{ = }{ \left(  \R^2  \right)^3
}
{  }{
}
{    }{
}
{   }{
}
}
{}{}{}
sebagai himpunan semua segitiga
\zusatzklammer {termasuk yang degenerat} {} {}
dengan mengidentifikasi segitiga yang mempunyai titik-titik sudut
\zusatzklammer {terurut} {} {}

\mavergleichskette
{\vergleichskette
{ P_1
}
{ = }{ \begin{pmatrix}  x_1  \\ y_1    \end{pmatrix}
}
{  }{
}
{    }{
}
{   }{
}
}
{}{}{,}

\mavergleichskette
{\vergleichskette
{ P_2
}
{ = }{ \begin{pmatrix}  x_2  \\ y_2    \end{pmatrix}
}
{  }{
}
{    }{
}
{   }{
}
}
{}{}{}
dan

\mavergleichskette
{\vergleichskette
{ P_3
}
{ = }{ \begin{pmatrix}  x_3  \\ y_3    \end{pmatrix}
}
{  }{
}
{    }{
}
{   }{
}
}
{}{}{,}
dengan tupel koordinat

\mathdisp {(x_1,y_1,x_2,y_2,x_3,y_3)} {  }

.
\aufzaehlungsechs{Tunjukkan bahwa himpunan segitiga dengan dua titik sudut berimpit merupakan subhimpunan tertutup $\R^6$
\zusatzklammer {jadi komplemennya merupakan himpunan terbuka dalam $\R^6$ yang kita namai $U$} {} {.}
}{Tunjukkan bahwa himpunan segitiga dengan ketiga titik sudut segaris merupakan subhimpunan tertutup $\R^6$
\zusatzklammer {jadi komplemennya, yang terdiri atas semua segitiga tak degenerat, merupakan himpunan terbuka dalam $\R^6$ yang kita namai $V$} {} {.}
}{Buat aturan fungsi yang mendeskripsikan pemetaan
\maabbdisp {F} {\R^6} {\R
} {}
yang memetakan suatu segitiga ke kelilingnya.
}{Tunjukkan bahwa fungsi $F$ dari bagian (3) terdiferensialkan secara kontinu pada himpunan $U$.
}{Hitung turunan parsial $F$ terhadap $x_1$ pada $U$.
}{Tunjukkan bahwa himpunan segitiga tak degenerat berkeliling $1$ membentuk
\definitionsverweis {submanifold tertutup}{}{}
dari $V$. Berapakah dimensinya?
}

}
{

\aufzaehlungsechs{Titik-titik sudut
\mathkor {} {P_1} {dan} {P_2} {}
berimpit jika dan hanya jika kedua koordinatnya sama, yakni jika

\mavergleichskette
{\vergleichskette
{ x_1
}
{ = }{ x_2
}
{  }{
}
{    }{
}
{   }{
}
}
{}{}{}
dan

\mavergleichskette
{\vergleichskette
{ y_1
}
{ = }{ y_2
}
{  }{
}
{    }{
}
{   }{
}
}
{}{}{}
berlaku. Syarat seperti ini mendefinisikan suatu himpunan bagian tertutup karena merupakan serat dari suatu pemetaan linear
\zusatzklammer {kontinu} {} {.}
Karena gabungan berhingga dari himpunan-himpunan bagian tertutup kembali tertutup, himpunan yang dimaksud bersifat tertutup.
}{Tinjau vektor-vektor penghubung

\mathdisp {P_2-P_1 = \begin{pmatrix}  x_2-x_1  \\ y_2-y_1   \end{pmatrix} \text{   und } P_3-P_1 = \begin{pmatrix}  x_3-x_1  \\ y_3-y_1   \end{pmatrix}} {  }

Ketiga titik kolinear jika dan hanya jika kedua vektor tersebut bergantung linear, dan hal ini berlaku jika dan hanya jika determinannya sama dengan $0$. Dengan demikian, syarat kolinearitasnya adalah

\mavergleichskettedisp
{\vergleichskette
{ H(x_1,x_2,x_3,y_1,y_2,y_3)
}
{ =} { (x_2-x_1) (y_3-y_1) - (x_3-x_1) (y_2-y_1)
}
{ =} { 0
}
{ } {
}
{ } {
}
}
{}{}{.}
Karena $H$ adalah pemetaan polinomial dan karenanya kontinu, serat di atas $0$ merupakan himpunan bagian tertutup.
}{Karena keliling merupakan jumlah dari ketiga sisi segitiga, berlaku

\mavergleichskettealigndrucklinks
{\vergleichskettealigndrucklinks
{ F(x_1,x_2,x_3,y_1,y_2,y_3)
}
{ =} { d(P_1,P_2) + d(P_1,P_3) +d(P_2,P_3)
}
{ =} { \sqrt{(x_2-x_1)^2 + (y_2-y_1)^2} +  \sqrt{(x_3-x_1)^2 + (y_3-y_1)^2} +  \sqrt{(x_3-x_2)^2 + (y_3-y_2)^2}
}
{ } {
}
{ } {
}
}
{}{}{.}
}{Untuk suatu titik di
\mathl{U}{,} setiap radikan, sebagai fungsi polinomial, terdiferensialkan secara kontinu dan bernilai positif ketat. Oleh karena itu, akar kuadratnya juga terdiferensialkan secara kontinu.
}{Turunan parsialnya adalah

\mavergleichskettealigndrucklinks
{\vergleichskettealigndrucklinks
{ { \frac{   \partial   F   }{  \partial   x_1  } }
}
{ =} { { \frac{   1 }{  2 } }   \left(  (  x_2-x_1  )^2 + (y_2-y_1)^2  \right)^{-1/2} (2x_1-2x_2)    + { \frac{   1 }{  2 } }   \left(  (  x_3-x_1  )^2 + (y_3-y_1)^2  \right)^{-1/2} (2x_1-2x_3)
}
{ =} {  \left(  (  x_2-x_1  )^2 + (y_2-y_1)^2  \right)^{-1/2} (x_1-x_2)    +   \left(  (  x_3-x_1  )^2 + (y_3-y_1)^2  \right)^{-1/2} (x_1-x_3)
}
{ } {
}
{ } {
}
}
{}{}{.}
}{Pada $U$, semua radikan bernilai positif, sehingga turunan parsial terhadap $x_1$ hanya dapat bernilai nol jika
\mathl{x_1=x_2=x_3}{.} Namun, segitiga yang memenuhi syarat ini bersifat kolinear. Ini berarti bahwa pada $V$ turunan parsial terhadap $x_1$ tidak pernah bernilai $0$; khususnya, gradien pada $V$ tidak pernah lenyap, sehingga bersifat reguler. Menurut
Satz 8.2 (Differentialgeometrie (Osnabrück 2023))
serat tersebut merupakan submanifold tertutup, dengan dimensi

\mavergleichskettedisp
{\vergleichskette
{ 6-1
}
{ =} { 5
}
{ } {
}
{ } {
}
{ } {
}
}
{}{}{.}
}

 }



\inputaufgabeklausurloesung
{4}

{

Misalkan
\mathkor {} {r} {dan} {R} {}
bilangan-bilangan real positif dengan

\mavergleichskette
{\vergleichskette
{ 0
}
{ < }{ r
}
{ < }{ R
}
{    }{
}
{   }{
}
}
{}{}{,}
dan tinjau torus
\zusatzklammer {tertanam} {} {}

\mavergleichskettedisp
{\vergleichskette
{ T
}
{ =} { \left\{ (x,y,z) \in \R^3  \mid   \left(  \sqrt{x^2+y^2} -R  \right)^2 +z^2  = r^2         \right\}
}
{ \subseteq} { \R^3
}
{ } {
}
{ } {
}
}
{}{}{}
dengan
\definitionsverweis {struktur Riemann terinduksi}{}{.}
Tunjukkan bahwa untuk setiap

\mavergleichskette
{\vergleichskette
{ \alpha
}
{ \in }{ [0,2 \pi [
}
{  }{
}
{    }{
}
{   }{
}
}
{}{}{}
pemetaan
\maabbeledisp {\Psi_\alpha} {\R^3} {\R^3
} { \left(  x   , \,   y  , \,  z                 \right) } { \left(  x \cos \alpha - y \sin  \alpha   , \,   x \sin \alpha + y \cos  \alpha  , \,  z                 \right)
} {,}
menginduksi suatu
\definitionsverweis {isometri}{}{}
dari $T$ ke dirinya sendiri.

}
{

Pada

\mavergleichskettedisp
{\vergleichskette
{ \R^3
}
{ =} { \R^2 \times \R
}
{ } {
}
{ } {
}
{ } {
}
}
{}{}{}
pemetaan tersebut merupakan produk suatu
\definitionsverweis {rotasi bidang}{}{}
dengan identitas. Oleh karena itu, pemetaan tersebut adalah
\definitionsverweis {isometri}{}{}
Euklides pada $\R^3$. Misalkan

\mavergleichskette
{\vergleichskette
{ (x,y,z)
}
{ \in }{  T
}
{  }{
}
{    }{
}
{   }{
}
}
{}{}{}
dan

\mavergleichskettedisp
{\vergleichskette
{ \Psi_\alpha(x,y,z)
}
{ =} { ( u,v,z)
}
{ } {
}
{ } {
}
{ } {
}
}
{}{}{.}
Maka

\mavergleichskettealigndrucklinks
{\vergleichskettealigndrucklinks
{ \left(  \sqrt{u^2+v^2} -R  \right)^2 +z^2
}
{ =} { \left(  \sqrt{ \left(  x \cos \alpha - y \sin  \alpha   \right)^2+ \left(  x \sin \alpha + y \cos  \alpha  \right)^2} -R  \right)^2 +z^2
}
{ =} { \left(  \sqrt{  x^2 \cos^{ 2  } \alpha + y^2 \sin^{ 2  }  \alpha  -2xy \cos \alpha \sin  \alpha    + x^2 \sin^{ 2  } \alpha + y^2 \cos^{ 2  }  \alpha  +2xy \cos \alpha \sin  \alpha     } -R  \right)^2 +z^2
}
{ =} { \left(  \sqrt{  x^2   + y^2 } -R  \right)^2 +z^2
}
{ =} { r^2
}
}
{}{}{,}
sehingga

\mavergleichskette
{\vergleichskette
{ \Psi_\alpha(x,y,z)
}
{ \in }{  T
}
{  }{
}
{    }{
}
{   }{
}
}
{}{}{.}
Jadi,
\maabbdisp {\Psi_\alpha} { T} {T
} {}
adalah suatu difeomorfisme dan isometri karena $T$ dilengkapi dengan struktur Riemann terinduksi.

 }



\inputaufgabeklausurloesung
{3}

{

Deskripsikan suatu model untuk bidang hiperbolik.

}
{  }



\inputaufgabeklausurloesung
{7}

{

Buktikan Theorema egregium.

}
{

Kita mendiferensialkan persamaan penentu pertama untuk simbol-simbol Christoffel, yaitu

\mavergleichskettedisp
{\vergleichskette
{ \partial_1 \partial_1 \varphi
}
{ =} { \Gamma^1_{11}  \partial_1 \varphi  + \Gamma^2_{11}  \partial_2 \varphi +h_{11} N
}
{ } {
}
{ } {
}
{ } {
}
}
{}{}{,}
terhadap variabel kedua, serta persamaan penentu kedua, yaitu

\mavergleichskettedisp
{\vergleichskette
{ \partial_1 \partial_2 \varphi
}
{ =} { \Gamma^1_{12}  \partial_1 \varphi + \Gamma^2_{12}  \partial_2  \varphi + h_{12} N
}
{ } {
}
{ } {
}
{ } {
}
}
{}{}{,}
terhadap variabel pertama. Berdasarkan teorema Schwarz, diperoleh

\mavergleichskettealigndrucklinks
{\vergleichskettealigndrucklinks
{ \left(  \partial_2 \Gamma^1_{11}  \right) \partial_1 \varphi +  \left(  \partial_2 \Gamma^2_{11}  \right)  \partial_2 \varphi+ \left(  \partial_2 h_{11}  \right)  N + \Gamma^1_{11} \partial_2 \partial_1 \varphi + \Gamma^2_{11}  \partial_2 \partial_2 \varphi  +h_{11} \partial_2 N
}
{ =} { \left(  \partial_1 \Gamma^1_{12}  \right)  \partial_1  \varphi+ \left(  \partial_1 \Gamma^2_{12}  \right)\partial_2 \varphi + \left(  \partial_1 h_{12}  \right)  N+ \Gamma^1_{12}  \partial_1 \partial_1 \varphi + \Gamma^2_{12}  \partial_1 \partial_2 \varphi + h_{12} \partial_1 N
}
{ } {
}
{ } {
}
{ } {
}
}
{}{}{.}
Selisih kedua ekspresi ini adalah $0$, dan kita menentukan bagian yang berkaitan dengan vektor basis $\partial_2 \varphi$. Untuk itu, kita menggunakan persamaan-persamaan penentu bagi simbol Christoffel dan
Lemma 19.3 (Differentialgeometrie (Osnabrück 2023))  (3)
, lalu memperoleh

\mavergleichskettedisphandlinks
{\vergleichskettedisphandlinks
{ 0
}
{ =} { \partial_2 \Gamma^2_{11} - \partial_1 \Gamma^2_{12} +\Gamma^1_{11} \Gamma^2_{12} - \Gamma^1_{12} \Gamma^2_{11} + \Gamma_{11}^2 \Gamma_{22}^2- \left(  \Gamma_{12}^2  \right)^2 + h_{11}  { \frac{   g_{12} h_{12} - g_{11} h_{22}   }{  g_{11}g_{22} -g_{12}^2 } } - h_{12}  { \frac{   g_{12} h_{11} - g_{11}h_{12}  }{  g_{11}g_{22} -g_{12}^2 } }
}
{ } {
}
{ } {
}
{ } {
}
}
{}{}{.}
Dengan

\mavergleichskettedisphandlinks
{\vergleichskettedisphandlinks
{ - g_{11} \left(  h_{11}h_{22} -h_{12}^2  \right)
}
{ =} { h_{11} \left(  g_{12} h_{12}- g_{11} h_{22}  \right) - h_{12} \left(  g_{12} h_{11} - g_{11}h_{12}  \right)
}
{ } {
}
{ } {
}
{ } {
}
}
{}{}{}
kita dapat mengganti kedua suku terakhir. Selanjutnya, menurut
Lemma 19.3 (Differentialgeometrie (Osnabrück 2023))  (4)
, diperoleh

\mavergleichskettealignhandlinks
{\vergleichskettealignhandlinks
{ g_{11} K
}
{ =} { g_{11} { \frac{   h_{11}h_{22} -h_{12}^2  }{ g_{11}g_{22} -g_{12}^2 } }
}
{ =} { \left(  \partial_2 \Gamma_{11}^2 - \partial_1 \Gamma_{12}^2 + \Gamma^1_{11} \Gamma^2_{12}  + \Gamma_{11}^2 \Gamma_{22}^2 - \Gamma_{12}^1 \Gamma_{11}^2-  \left(  \Gamma_{12}^2  \right)^2  \right)
}
{ } {
}
{ } {
}
}
{}
{}{.}

 }



\inputaufgabeklausurloesung
{6 (1+2+2+1)}

{

Kita tinjau
\definitionsverweis {bentuk diferensial}{}{}

\mavergleichskettedisp
{\vergleichskette
{ \omega
}
{ =} { ydx+zdy +xdz
}
{ } {
}
{ } {
}
{ } {
}
}
{}{}{}
pada $\R^3$ dan pemetaan
\maabbeledisp {\varphi} {\R^2} {\R^3
} { \left(  u   , \,  v                   \right) } { \left( u^2  , \,  v^2  , \,  uv                 \right)
} {.}
\aufzaehlungvier{Hitung turunan eksterior $\omega$.
}{Hitung tarikan balik $\varphi^* \omega$ dari $\omega$ di bawah $\varphi$.
}{Hitung turunan eksterior $\varphi^* \omega$ pada $\R^2$.
}{Hitung tarikan balik $\varphi^* d \omega$ dari $d\omega$ di bawah $\varphi$ secara independen dari bagian (3).
}

}
{

\aufzaehlungvier{Berlaku

\mavergleichskettedisp
{\vergleichskette
{ d \omega
}
{ =} { d \left(  ydx+zdy +xdz  \right)
}
{ =} { dy \wedge dx +dz \wedge dy +dx \wedge dz
}
{ =} { - dx \wedge dy -dy \wedge dz +dx \wedge dz
}
{ } {
}
}
{}{}{.}
}{Berlaku

\mavergleichskettealign
{\vergleichskettealign
{  \varphi^* \omega
}
{ =} { v^2 d u^2 + uvdv^2+ u^2duv
}
{ =} { 2v^2udu +2uv^2 dv + u^2(udv +vdu)
}
{ =} { \left(  2uv^2 +u^2v  \right) du + \left(  2uv^2 +u^3  \right) dv
}
{ } {
}
}
{}
{}{.}
}{Berlaku

\mavergleichskettealign
{\vergleichskettealign
{  d \left(   \left(  2uv^2 +u^2v  \right) du + \left(  2uv^2 +u^3  \right) dv     \right)
}
{ =} {  d \left(  2uv^2 +u^2v  \right) du  +  d\left(  2uv^2 +u^3  \right) dv
}
{ =} { \left(  4uv  + u^2      \right)  dv \wedge du  +   \left(  2v^2   +3u^2  \right)   du \wedge dv
}
{ =} { \left(  -4uv +  2v^2   +2u^2  \right)   du \wedge dv
}
{ } {
}
}
{}
{}{.}
}{Tarikan balik $\varphi^* d \omega$ adalah

\mavergleichskettealign
{\vergleichskettealign
{ \varphi^* \left(  - dx \wedge dy -dy \wedge dz +dx \wedge dz   \right)
}
{ =} { - d u^2 \wedge dv^2 -d v^2 \wedge duv +d u^2 \wedge duv
}
{ =} { -4uv  du \wedge dv -2v^2 dv \wedge du + 2   u^2 du \wedge dv
}
{ =} { \left(  -4uv+2v^2 + 2 u^2  \right)  du \wedge dv
}
{ } {
}
}
{}
{}{.}
}

 }



\inputaufgabeklausurloesung
{3}

{

Misalkan $M$ suatu
\definitionsverweis {manifold diferensiabel}{}{}
dan $\omega$ suatu
\definitionsverweis {bentuk diferensial}{}{}
yang
\definitionsverweis {eksak}{}{}
pada $M$. Misalkan $S$ suatu
\definitionsverweis {manifold diferensiabel}{}{}
yang
\definitionsverweis {kompak}{}{}
dan
\definitionsverweis {berorientasi}{}{}
\zusatzklammer {tanpa batas} {} {}
dengan
\definitionsverweis {basis terhitung bagi topologinya}{}{}
dan misalkan
\maabbdisp {\varphi} { S } { M
} {}
suatu
\definitionsverweis {pemetaan terdiferensialkan secara kontinu}{}{.}
Tunjukkan bahwa

\mavergleichskettedisp
{\vergleichskette
{
\int_{  S  }   \varphi^*\omega
}
{ =} { 0
}
{ } {
}
{ } {
}
{ } {
}
}
{}{}{.}

}
{

Karena menurut
Satz 20.4 (Differentialgeometrie (Osnabrück 2023))  (5)
operasi tarikan balik bentuk kompatibel dengan turunan eksterior,
\mathl{\varphi^*\omega}{} juga eksak. Misalkan $\tau$ adalah suatu bentuk diferensiabel pada $S$ dengan

\mavergleichskettedisp
{\vergleichskette
{d \tau
}
{ =} { \varphi^* \omega
}
{ } {
}
{ } {
}
{ } {
}
}
{}{}{.}
Menurut
Satz von Stokes,
berlaku

\mavergleichskettedisp
{\vergleichskette
{ \int_S \varphi^* \omega
}
{ =} { \int_S d \tau
}
{ =} { \int_{ \partial S } \tau
}
{ =} { \int_\emptyset \tau
}
{ =} { 0
}
}
{}{}{.}

 }



\inputaufgabeklausurloesung
{6}

{

Buktikan teorema tentang retraksi ke batas pada manifold.

}
{

Batas
\mathl{\partial M}{}, menurut
Satz 21.8 (Differentialgeometrie (Osnabrück 2023))
, merupakan suatu
\definitionsverweis {manifold diferensiabel}{}{}
\definitionsverweis {berorientasi}{}{}
\zusatzklammer {tanpa batas} {} {.}
Oleh karena itu, menurut
Satz 22.11 (Differentialgeometrie (Osnabrück 2023))
, terdapat suatu
\definitionsverweis {bentuk volume positif}{}{}
\definitionsverweis {yang terdiferensialkan secara kontinu}{}{}
$\tau$ pada
\mathl{\partial M}{.} Kita mempunyai
\mathl{\int_{  \partial M  }  \tau >0}{.}
\definitionsverweis {Turunan eksterior}{}{}
dari bentuk volume $\tau$ adalah $0$. Andaikan terdapat suatu
\definitionsverweis {pemetaan yang terdiferensialkan secara kontinu}{}{}
\maabbdisp {\varphi} { M } { \partial M
} {}
dengan
\mathl{\varphi \vert_{\partial M} = \operatorname{Id}_{\partial M}}{.} Maka
\definitionsverweis {bentuk tarikan balik}{}{}

\mathl{\varphi^* \tau}{} adalah
$(n-1)$-\definitionsverweis {bentuk diferensial}{}{}
pada $M$ yang pembatasannya pada batas berimpit dengan $\tau$. Dengan menggunakan
Satz 23.2 (Differentialgeometrie (Osnabrück 2023))
dan
Satz 20.4 (Differentialgeometrie (Osnabrück 2023))  (5)
kita memperoleh

\mavergleichskettealign
{\vergleichskettealign
{
\int_{  \partial M  }  \tau
}
{ =} {
\int_{  \partial  M  }  \varphi^* \tau
}
{ =} {
\int_{   M  }  d(\varphi^* \tau)
}
{ =} {
\int_{   M  }  \varphi^* (d\tau)
}
{ =} {
\int_{   M  }  \varphi^* (0)
}
}
{
\vergleichskettefortsetzungalign
{ =} { 0
}
{ } {}
{ } {}
{ } {}
}
{}{.}
Ini merupakan suatu kontradiksi.

 }
